% CREACIÓN DEL DOCUMENTO
\documentclass[
	spanish, % Idioma: spanish, english, etc.
	letterpaper, oneside
]{article}

% INFORMACIÓN DEL DOCUMENTO
\def\documenttitle {Actividad 1: Entradas y salidas de un programa}
\def\documentsubtitle {}
\def\documentsubject {Entradas y salidas de un programa}
\def\documentdate {\today}

\def\documentauthor {Emilio Abraham Castillo Urdiales}
\def\coursename {Metodología de la Programación}
\def\coursecode {Enero -Junio 2026}

\def\universityname {Universidad Autónoma de Nuevo Léon}
\def\universityfaculty {Facultad de Ingeniería Mecánica y Eléctrica}
\def\universitydepartment {Ingeniería en Administración y Sistemas}
\def\universitydepartmentimage {departamentos/fcfm}
\def\universitylocation {San Nicolás de los Garza, Nuevo León, México}

% IMPORTACIÓN DEL TEMPLATE
\input{template}

% INICIO DE PÁGINAS
\begin{document}

% CONFIGURACIÓN DE PÁGINA Y ENCABEZADOS
\templatePagecfg

% CONFIGURACIONES FINALES
\templateFinalcfg

% ======================= Descripción de la actividad ================
%Cada programa deberá ser entregado en la actividad programada en Teams como "Actividad fundamental 1"
%Se deberán entregar 7 programas (.psc) Generados por PSEINT
%Además de esto, deberán entregar un PDF con todos los diagramas, pseudocodigos y ejecuciones de sus programas especificando cual programa es cada uno.
%
%Ejemplo:
%Programa 1:
%*Foto de pseudocódigo (Screenshot)
%*Foto de diagrama (Screenshot)
%*Foto de la ejecución del programa (Screenshot)
%Programa 2:
%*Foto de pseudocodigo (Screenshot)
%*Foto de diagrama (Screenshot)
%*Foto de la ejecución del programa (Screenshot)
%Así con los 7 programas
%RESUMEN: se entregarán:
%7 archivos .psc (no comprimidos, si se entregan comprimidos se bajarán puntos)
%1 PDF (no .Docx ni .Pptx)
%||||||||||||||||||||||||||||||||||||||||||||||||||||||||||||||||||||||||||||||||||||||||||||||||||||||||||||||||||||||||||||||
%Programa 1: Suma de dos números y mostrar el resultado
%Programa 2: Suma, Resta, Multiplicación y división de dos números (mostrar los 4 resultados)
%Programa 3: Programa el cual calcule la energía potencial de un cuerpo en caída libre. (Epot = mgh)
%
%Programa 4: Programa Que calcule e imprima el volumen de una esfera vol $=4 / 3 \Pi r^3$
%Programa 5: La compañía de la luz desea generar el recibo de un cliente a partir de los siguientes datos:
%Nombre del cliente
%Consumo en kilowatts
%Tarifa del kilowatt.
%El recibo solo deberá contener el nombre y pago del cliente.
%Programa 6: Una persona compró un terreno en un país. La extensión del terreno está especificada en acres. Introducir como dato la extensión del terreno en acres, calcular e imprimir la extensión del mismo en hectáreas.
%- Un acre $=4047 \mathrm{~m} 2$.
%- Una hectárea $=10000 \mathrm{~m} 2$.
%
%Programa 7: Introducir como datos los valores de la longitud y el peso de un objeto expresado en pies y libras, imprimir los datos de este objeto pero expresado en metros y kilogramos.
%- Un pie $=0.3048$ Mts.
%- Una libra $=0.45359 \mathrm{Kgs}$.

% ======================= Ejemplos de pseudocódigo con Pseint =======================
%// este es el ejemplo más simple de esta ayuda, 
%// toma dos numeros, los suma y muestra el resultado
%

%Algoritmo Suma
%
%
%// para cargar un dato, se le muestra un mensaje al usuario
%// con la instrucción Escribir, y luego se lee el dato en 
%// una variable (A para el primero, B para el segundo) con 
%// la instrucción Leer
%
%Escribir "Ingrese el primer numero:"
%Leer A
%
%Escribir "Ingrese el segundo numero:"
%Leer B
%
%
%// ahora se calcula la suma y se guarda el resultado en la
%// variable C mediante la asignación (<-)
%
%C <- A+B
%
%
%// finalmente, se muestra el resultado, precedido de un 
%// mensaje para avisar al usuario, todo en una sola
%// instrucción Escribir
%
%Escribir "El resultado es: ",C
%
%FinAlgoritmo

%// Este ejemplo muestra el uso de expresiones, operadores y funciones matematicas
%
%Algoritmo Matematicas
%
%
%EligeSalir<-Falso
%Escribir 'Ingresar un número:'
%Leer N
%Repetir
%Escribir ' '
%Escribir 'Presione una tecla para continuar'
%Esperar Tecla
%Limpiar Pantalla
%Escribir 'Elija una opción:'
%Escribir '  1 - Seno, coseno, arcotangente'
%Escribir '  2 - Lograritmo natural, función exponencial'
%Escribir '  3 - Truncar, redondear'
%Escribir '  4 - Raíz cuadrada'
%Escribir '  5 - Valor absoluto'
%Escribir '  6 - Separar parte entera y decimal'
%Escribir '  7 - Hallar factorial'
%Escribir '  8 - Averiguar si es primo'
%Escribir '  9 - Ingresar otro número'
%Escribir '  0 - Salir'
%Escribir ' '
%Leer eleccion
%Segun eleccion Hacer
%1: 
%Escribir 'Seno:',Sen(N)
%Escribir 'Coseno:',Cos(N)
%Escribir 'Arcotangente:',Atan(N)
%2: 
%Si N<=0
%Entonces Escribir 'El numero debe ser mayor a cero!'
%SiNo
%Escribir 'Logaritmo natural: ',ln(N)
%Escribir 'Función exponencial: ',exp(N)
%FinSi
%3: 
%Escribir 'Truncar: ',trunc(N)
%Escribir 'Redondear: ',redon(N)
%4: Escribir 'Raiz Cuad.: ',rc(N)
%5: Escribir 'Valor Abs.: ',abs(N)
%6: 
%Escribir 'Parte Entera: ',Trunc(n)
%Escribir 'Parte Decimal: ',n-Trunc(n)
%7: 
%Si N<>Trunc(N)
%Entonces
%Escribir 'El numero debe ser entero!'
%SiNo
%Si abs(N)>50
%Entonces Escribir 'Resultado muy grande!'
%SiNo
%r<-1; f<-1
%Mientras f<=abs(N) Hacer
%Si N<0 
%Entonces r<-(-f)*r
%SiNo
%r<-f*r
%FinSi
%f<-f+1
%FinMientras
%Escribir 'Factorial:',r
%FinSi
%FinSi
%8: 
%Si N<>Trunc(N) Entonces
%Escribir 'El numero debe ser entero!'
%SiNo
%Si N<0  entonces 
%Nu<-N*(-1)
%SiNo
%Nu<-N
%FinSi
%Si N mod 2 = 0 Entonces 
%Escribir 'Numero Primo:',Nu=2
%Si Nu<>2 Entonces
%Escribir N,'=2x',N/2
%FinSi
%SiNo
%EsPrimo<-Nu<>1 
%Nu<-RC(Nu)
%Divisor<-3
%Mientras Divisor<=Raiz(Nu) & EsPrimo Hacer
%Si N mod Divisor = 0 Entonces 
%EsPrimo<-Falso
%Sino
%Divisor<-Divisor+2
%FinSi
%FinMientras
%Escribir 'Numero primo:',EsPrimo
%Si N>1 & ~ EsPrimo Entonces 
%Escribir N,'=',Divisor,'x',N/Divisor
%FinSi
%FinSi
%FinSi
%9:
%Escribir 'Ingrese un número:'
%Leer N
%0: EligeSalir<-Verdadero
%De Otro Modo:
%Escribir 'eleccionción  no válida!'
%FinSegun
%Hasta que EligeSalir
%FinAlgoritmo


% ======== Ejemplo de insertar bloque de código pseudocódigo en el documento =======================
%\begin{sourcecode}{pseudocodecolor}{Algoritmo con fondo de color.}
%	input: int N, int D
%	output: int
%	begin
%	res $\gets$ 0
%	while N $\geq$ D 
%	N $\gets$ N - D
%	res $\gets$ res + 1      
%	end
%	return res
%	end    
%\end{sourcecode}

% ======================= INICIO DEL DOCUMENTO =======================

% Título y nombre del autor
\inserttitle

% Resumen o Abstract

\begin{abstract}
	
	Se realizan 7 programas en PSEINT, cada uno con su respectivo pseudocódigo, diagrama de flujo y ejecución. Los programas son los siguientes:
	\begin{enumerate}
		\item Suma de dos números.
		\item Operaciones básicas (suma, resta, multiplicación, división).
		\item Energía potencial en caída libre ($E_{pot} = mgh$).
		\item Volumen de una esfera ($V = \frac{4}{3}\pi r^3$).
		\item Generación de recibo de luz.
		\item Conversión de acres a hectáreas.
		\item Conversión de pies y libras a metros y kilogramos.
	\end{enumerate}

\end{abstract}

%\begin{multicols}{2}

\section{Suma de dos números}

El programa solicita al usuario ingresar dos números, realiza la suma de ambos y muestra el resultado en pantalla. El pseudocódigo, diagrama de flujo y ejecución se presentan a continuación.		
	
\begin{sourcecode}{pseudocodecolor}{Pseudocódigo del programa de suma de dos números.}
Algoritmo SumaDosNumeros
    Escribir "Ingrese el primer número:"
    // Escribe el primer número en pantalla y lo lee desde el teclado
    Leer numero1
    Escribir "Ingrese el segundo número:"
    // Escribe el segundo número en pantalla y lo lee desde el teclado
    Leer numero2
    // Realiza la suma de los dos números y almacena el resultado en la variable suma
    suma <- numero1 + numero2
    // Escribe el resultado de la suma en pantalla
    Escribir "El resultado de la suma es: ", suma
FinAlgoritmo
\end{sourcecode}

\insertimage[\label{img:diagramaSumaDosNumeros}]{wVTA6MDTxNuGsL_CoY_f9eVqoQwKmprMQ0_UBUYD8Mc_original_fullsize}{width=0.75\textwidth}{Diagrama de flujo del programa de suma de dos números.}

\insertimage[\label{img:ejecucionSumaDosNumeros}]{QRmCLQYLN09bKG1r-KeCjKUBVgviZE3uglSzaakOGUE_original_fullsize}{width=0.45\textwidth}{Ejecución del programa de suma de dos números.}

\section{Suma, Resta, Multiplicación y División de dos números}

El programa solicita al usuario ingresar dos números, realiza las operaciones de suma, resta, multiplicación y división, y muestra los resultados en pantalla. El pseudocódigo, diagrama de flujo y ejecución se presentan a continuación.

\begin{sourcecode}{pseudocodecolor}{Pseudocódigo del programa de operaciones básicas.}
Algoritmo OperacionesBasicas
    Escribir "Ingrese el primer número:"
    // Lee el primer valor ingresado por el usuario
    Leer numero1
    Escribir "Ingrese el segundo número:"
    // Lee el segundo valor ingresado por el usuario
    Leer numero2
    // Calcula la suma de ambos números
    suma <- numero1 + numero2
    // Calcula la resta del primer número menos el segundo
    resta <- numero1 - numero2
    // Calcula la multiplicación de ambos números
    multiplicacion <- numero1 * numero2
    // Valida que el divisor sea diferente de cero antes de dividir
    Si numero2 <> 0 Entonces
        // Realiza la división cuando es válida
        division <- numero1 / numero2
    Sino
        // Muestra un mensaje de error si el divisor es cero
        Escribir "No se puede dividir por cero"
    FinSi
    // Muestra todos los resultados calculados
    Escribir "Suma: ", suma
    Escribir "Resta: ", resta
    Escribir "Multiplicación: ", multiplicacion
    Escribir "División: ", division
FinAlgoritmo
\end{sourcecode}

\insertimage[\label{img:diagramaOperacionesBasicas}]{VlDw_pMjVS0DTo9f96Sqe7TrbfeD021Rf_F-i0c3AG8_original_fullsize}{width=0.75\textwidth}{Diagrama de flujo del programa de operaciones básicas.}

\insertimage[\label{img:ejecucionOperacionesBasicas}]{g4BahtJicLW6VnQD7Phe7xrrqIutu_A_lt0i_BFeH5c_original_fullsize}{width=0.45\textwidth}{Ejecución del programa de operaciones básicas.}

\section{Energía potencial de un cuerpo en caída libre}

El programa calcula la energía potencial de un cuerpo en caída libre utilizando la fórmula $E_{pot} = mgh$, donde $m$ es la masa del cuerpo, $g$ es la aceleración debido a la gravedad (9.81 m/s²) y $h$ es la altura desde la cual cae el cuerpo. El pseudocódigo, diagrama de flujo y ejecución se presentan a continuación.

\begin{sourcecode}{pseudocodecolor}{Pseudocódigo del programa de energía potencial.}
Algoritmo EnergiaPotencial
	Escribir "Ingrese la masa del cuerpo (en kg):"
	// Lee la masa del cuerpo ingresada por el usuario
	Leer masa
	Escribir "Ingrese la altura desde la cual cae el cuerpo (en metros):"
	// Lee la altura desde la cual cae el cuerpo ingresada por el usuario
	Leer altura
	// Define la constante de gravedad
	gravedad <- 9.81
	// Calcula la energía potencial utilizando la fórmula Epot = mgh
	energia_potencial <- masa * gravedad * altura
	// Muestra el resultado de la energía potencial en pantalla
	Escribir "La energía potencial del cuerpo es: ", energia_potencial, " Joules"
FinAlgoritmo
\end{sourcecode}

\insertimage[\label{img:diagramaEnergiaPotencial}]{6_rVw212epf2vyUVNSIDrqbcmMXGcnvgFx9RHssXgCg_original_fullsize}{width=0.75\textwidth}{Diagrama de flujo del programa de energía potencial.}

\insertimage[\label{img:ejecucionEnergiaPotencial}]{ZEAn9IWjIodIKlQt6pWWjhDDPMShedOIDWDW5c1VS5A_original_fullsize}{width=0.45\textwidth}{Ejecución del programa de energía potencial.}

\section{Volumen de una esfera}

El programa calcula el volumen de una esfera utilizando la fórmula $V = \frac{4}{3}\pi r^3$, donde $r$ es el radio de la esfera. El pseudocódigo, diagrama de flujo y ejecución se presentan a continuación.

\begin{sourcecode}{pseudocodecolor}{Pseudocódigo del programa de volumen de una esfera.}
Algoritmo calculoVolumenEsfera
    Escribir "Ingrese el radio de la esfera (en metros):"
    // Lee el radio de la esfera ingresado por el usuario
    Leer radio
    // Calcula el volumen de la esfera utilizando la fórmula vol = (4/3) * π * r^3
    volumen <- (4/3) * 3.14159 * (radio^3)
    // Muestra el resultado del volumen calculado
    Escribir "El volumen de la esfera es: ", volumen, " metros cúbicos"
FinAlgoritmo
\end{sourcecode}

\insertimage[\label{img:diagramaVolumenEsfera}]{0IH7toh9jWq0SrzXVHQ4bdC4FBfhl3t72cQ7H-4tj_U_original_fullsize}{width=0.75\textwidth}{Diagrama de flujo del programa de volumen de una esfera.}

\insertimage[\label{img:ejecucionVolumenEsfera}]{-57AH2iXf95ibZNPnJ48dlz0lr8et5lp83lDMfH6MaU_original_fullsize}{width=0.45\textwidth}{Ejecución del programa de volumen de una esfera.}

\section{Recibo de luz}

El programa genera un recibo de luz a partir de los datos ingresados por el usuario: nombre del cliente, consumo en kilowatts y tarifa del kilowatt. El recibo muestra el nombre del cliente y el pago total. El pseudocódigo, diagrama de flujo y ejecución se presentan a continuación.

\begin{sourcecode}{pseudocodecolor}{Pseudocódigo del programa de recibo de luz.}
Algoritmo generarReciboLuz
    Escribir "Ingrese el nombre del cliente:"
    // Lee el nombre del cliente ingresado por el usuario
    Leer nombreCliente
    Escribir "Ingrese el consumo en kilowatts:"
    // Lee el consumo en kilowatts ingresado por el usuario
    Leer consumoKilowatts
    Escribir "Ingrese la tarifa del kilowatt:"
    // Lee la tarifa del kilowatt ingresada por el usuario
    Leer tarifaKilowatt
    // Calcula el pago total multiplicando el consumo por la tarifa
    pagoTotal <- consumoKilowatts * tarifaKilowatt
    // Muestra el recibo con el nombre del cliente y el pago total
    Escribir "Recibo de la compañía de la luz"
    Escribir "Cliente: ", nombreCliente
    Escribir "Pago total: ", pagoTotal
FinAlgoritmo
\end{sourcecode}

\insertimage[\label{img:diagramaReciboLuz}]{VPCyOpBiWVDx3o0B4pRpBT5xGlO0iJVtVdFnh1Iecbk_original_fullsize}{width=0.75\textwidth}{Diagrama de flujo del programa de recibo de luz.}

\insertimage[\label{img:ejecucionReciboLuz}]{0X7n2SguGBr7fJO7MA9w6AawY7YvsEjsei3EmUrH7-0_original_fullsize}{width=0.45\textwidth}{Ejecución del programa de recibo de luz.}

\section{Conversión de acres a hectáreas}

El programa convierte la extensión de un terreno de acres a hectáreas utilizando las equivalencias: 1 acre = 4047 m² y 1 hectárea = 10000 m². El programa solicita al usuario ingresar la extensión del terreno en acres, realiza la conversión y muestra el resultado en pantalla en hectáreas. El pseudocódigo, diagrama de flujo y ejecución se presentan a continuación.

\begin{sourcecode}{pseudocodecolor}{Pseudocódigo del programa de conversión de acres a hectáreas.}
Algoritmo conversionAcresHectareas
    Escribir "Ingrese la extensión del terreno en acres:"
    // Lee la extensión del terreno en acres ingresada por el usuario
    Leer extensionAcres
    // Convierte la extensión de acres a hectáreas utilizando la equivalencia 1 acre = 0.4047 hectáreas
    extensionHectareas <- extensionAcres * 0.4047
    // Muestra el resultado de la conversión en pantalla
    Escribir "La extensión del terreno en hectáreas es: ", extensionHectareas, " ha"
FinAlgoritmo
\end{sourcecode}

\insertimage[\label{img:diagramaConversionAcresHectareas}]{TlppQIz-eW4O78lhtq5lW5ghqqjOnd57OmoTqxjdPgU_original_fullsize}{width=0.75\textwidth}{Diagrama de flujo del programa de conversión de acres a hectáreas.}

\insertimage[\label{img:ejecucionConversionAcresHectareas}]{6ixxdOeiDU_dlse8QuXMQ8Es0o8gXuFspW5Ug76Kx04_original_fullsize}{width=0.45\textwidth}{Ejecución del programa de conversión de acres a hectáreas.}
%
\section{Conversión de pies y libras a metros y kilogramos}

El programa solicita al usuario ingresar la longitud en pies y el peso en libras, realiza la conversión a metros y kilogramos utilizando las equivalencias dadas, y muestra los resultados en pantalla. El pseudocódigo, diagrama de flujo y ejecución se presentan a continuación.

\begin{sourcecode}{pseudocodecolor}{Pseudocódigo del programa de conversión de pies y libras a metros y kilogramos.}
Algoritmo conversionPiesLibras
    Escribir "Ingrese la longitud en pies:"
    // Lee la longitud en pies ingresada por el usuario
    Leer longitudPies
    Escribir "Ingrese el peso en libras:"
    // Lee el peso en libras ingresado por el usuario
    Leer pesoLibras
    // Convierte la longitud de pies a metros utilizando la equivalencia 1 pie = 0.3048 metros
    longitudMetros <- longitudPies * 0.3048
    // Convierte el peso de libras a kilogramos utilizando la equivalencia 1 libra = 0.45359 kilogramos
    pesoKilogramos <- pesoLibras * 0.45359
    // Muestra los resultados de la conversión en pantalla
    Escribir "La longitud en metros es: ", longitudMetros, " m"
    Escribir "El peso en kilogramos es: ", pesoKilogramos, " kg"
FinAlgoritmo
\end{sourcecode}

\insertimage[\label{img:diagramaConversionPiesLibras}]{WQhtTmkyyyRaY4ecjoMBHaPLd7plI45D--nv3TSS3Tg_original_fullsize}{width=0.75\textwidth}{Diagrama de flujo del programa de conversión de pies y libras a metros y kilogramos.}

\insertimage[\label{img:ejecucionConversionPiesLibras}]{6TEUI5N0AJjGVBWvXW2EhrcY4HKP0jGRqrtdIqXDXa4_original_fullsize}{width=0.45\textwidth}{Ejecución del programa de conversión de pies y libras a metros y kilogramos.}

%\end{multicols}

% FIN DEL DOCUMENTO
\end{document}